% ============================================================
\subsection{Comparative Persistence and Gauge-Curvature Realization}
\label{sec:flagship-gauge}
% ============================================================

The gauge lane is the realization of \emph{comparative persistence}: only the
local comparison class is physically visible, and coherent exchange persists as
visible curvature and current. Exactness, local gauge covariance, and
inherited minimal-order closure then determine the physical
curvature/exchange law class on that carrier. The final carrier-level
identification of the canonical representative is completed below by local
endpoint rigidity on the selected one-form carrier.

Theorem~\ref{thm:temporal-realization-forbids-primitive-internal-creation}
identifies admissible one-form sources on that cut as visible exchange
currents rather than primitive creations.

The relevant compatibility identity on this carrier is nilpotent exactness:
closed curvature is inherited from the complex law \(d^2=0\) already built in
the earlier structural chapters together with the continuum bridge. The
spacetime slab \(M=J\times\Sigma\) introduced below is a
coordinate model of the realized temporal cut, with \(J\) carrying the
realized causal order.

\paragraph{Carrier status.}
The one-form carrier is likewise not being postulated independently here. By
Theorem~\ref{thm:three-forcing-bridges-from-bedrock-to-primitive-physics}, the
primitive direct lanes are already forced on the selected bundle, and the
gauge lane is the canonical realization of the \(3\)-channel. The role of the
present flagship section is to identify the gauge-compatible local datum and
its canonical representative, not to add a new sector law beside the temporal
realization principle.

\subsubsection{Gauge-Compatible Realization}

\begin{definition}[Gauge-compatible one-form realization]
\label{def:gauge-compatible-one-form-realization}\lean{GaugeLane.OneFormRealization}
Let \(M=J\times\Sigma\) be an oriented pseudo-Riemannian spacetime slab with
exterior derivative \(d\) and codifferential \(\delta\) induced by the Hodge
star of the chosen metric. A
\emph{gauge-compatible one-form realization} consists of:
\begin{enumerate}[label=(\roman*),itemsep=0.2em]
\item a visible carrier of local one-form classes \([A]\in\Omega^1(M)/d\Omega^0(M)\);
\item an exchange one-form \(J\in\Omega^1(M)\) satisfying
      \[
      \delta J=0;
      \]
\item the compatibility package given by exactness
      \[
      d^2=0,
      \]
      local gauge covariance under \(A\mapsto A+d\chi\), and continuum
      equivalence modulo local gauge;
\item the realized class is transport-order disciplined on the selected
      one-form carrier, and its distinguished admissible minimal-order class
      is local order \(2\) in \(A\), hence local order \(1\) in the
      transported curvature \(F:=dA\).
\end{enumerate}
\end{definition}

The codifferential \(\delta\) therefore depends on the metric/Hodge data on
\(M\); it is not determined by orientation alone.

\subsubsection{Least-Motion Gauge Carrier}

\begin{theorem}[Least-motion gauge observer]
\label{thm:least-motion-gauge-observer}\lean{GaugeLane.leastMotionObserver}
Assume the hypotheses of
Definition~\ref{def:gauge-compatible-one-form-realization}. Suppose:
\begin{enumerate}[label=(\roman*),itemsep=0.2em]
\item the realized family satisfies the hypotheses of
      Corollary~\ref{cor:intrinsic-package-yields-gauge-compatible-local-datum},
      so the selected law already carries a unique gauge-compatible local
      internal datum on the relevant selected channels;
\item the one-form carrier is faithful, closure-admissible, and stable under
      admissible promotion for the realized family;
\item no proper visible subcarrier still supports gauge-covariant,
      exactness-compatible exchange response together with faithful
      reconstruction and stable promotion at the observed scale;
\item any characteristic observer datum satisfying \emph{(i)}--\emph{(iii)} is
      horizon-preserving equivalent to the one generated by the selected
      one-form carrier.
\end{enumerate}
Then the selected one-form carrier is the unique least-motion characteristic
observer for the realized family.
\end{theorem}

(Proof in Appendix~\ref{sec:appendix-carrier}.)

\begin{definition}[Gauge-compatible curvature/exchange datum]
\label{def:gauge-compatible-curvature-source-datum}\lean{GaugeLane.CurvatureSourceDatum}
Assume the hypotheses of
Theorem~\ref{thm:least-motion-gauge-observer}. The
\emph{gauge-compatible curvature/exchange datum} on the selected cut is the
triple
\[
\mathfrak M_{\mathrm g}:=([A],F,J),
\qquad
F:=dA,
\]
where \([A]\) is the selected local gauge class on the one-form carrier and
\(J\) is the selected visible exchange current on that same cut.
\end{definition}

\subsubsection{Induced Curvature/Source Law}

\begin{theorem}[Induced gauge-compatible law on the selected cut]
\label{thm:induced-gauge-compatible-law}\lean{GaugeLane.inducedLaw}
Assume the hypotheses of
Theorem~\ref{thm:least-motion-gauge-observer}. Let
\(\mathfrak M_{\mathrm g}=([A],F,J)\) be the selected
gauge-compatible curvature/exchange datum of
Definition~\ref{def:gauge-compatible-curvature-source-datum}. Then there are a
unique gauge-compatible local response operator
\[
\mathcal M_{\mathrm g}:\Omega^1(M)/d\Omega^0(M)\to\Omega^1(M)
\]
of minimal local order and a local gauge-invariant remainder
\(\mathsf R_{\mathrm g}\in\Omega^1(M)\) such that
\[
dF=0,
\qquad
\mathcal M_{\mathrm g}([A])=J+\mathsf R_{\mathrm g}.
\]
Here \(J\) is the selected visible exchange current, and
Theorem~\ref{thm:temporal-realization-forbids-primitive-internal-creation}
identifies \(\mathsf R_{\mathrm g}\) as the remaining local gauge-invariant
exchange remainder on the selected cut.
If the selected cut is lossless at leading order, so that
\(\mathsf R_{\mathrm g}=0\), this is the forced gauge-compatible law on the
selected cut.
\end{theorem}

(Proof in Appendix~\ref{sec:appendix-carrier}.)

Thus the gauge carrier is already reduced to a local representative problem on
the selected cut: identify the canonical gauge-compatible representative
inside the forced curvature/exchange closure class. The local collapse below
first identifies the intrinsic selected-cut curvature-adjoint representative;
the recognizable Maxwell notation is then a downstream Hodge-side readout of
that already forced representative.

\subsubsection{Endpoint Identification}

In differential-form notation, the recognized endpoint law is the classical
Maxwell curvature/source system \cite{Maxwell1865}. The final response
operator is now read through a final local rigidity step. Exactness and local
gauge covariance already reduce the response from the gauge class \([A]\) to
the curvature \(F=dA\). The final identification is the unique
pairing-compatible first-order curvature representative on the selected cut.
When the generated one-form carrier is read on an energy-orthogonal realized
cut with its Hodge-side duality structure fixed, that representative is the
codifferential-curvature response.

\begin{definition}[Gauge local endpoint-rigidity package]
\label{def:gauge-local-endpoint-rigidity-package}\lean{GaugeLane.EndpointCollapse}
Assume the hypotheses of Theorem~\ref{thm:induced-gauge-compatible-law}. A
\emph{gauge local endpoint-rigidity package} on the selected one-form carrier
consists of:
\begin{enumerate}[label=(\roman*),itemsep=0.2em]
\item the realized family satisfies a canonical selected Schur bridge on that
      same selected cut;
\item the minimal pairing-symmetric response family on the selected one-form
      carrier is represented levelwise by a refinement-compatible local
      gauge-covariant stencil family of minimal local order \(2\) in \(A\),
      hence local order \(1\) in the curvature \(F=dA\);
\item the distinguished minimal-order response family on the selected one-form
      carrier carries a refinement-coherent reconstruction datum on its dense
      test domain whose exact core realization is bounded.
\end{enumerate}
\end{definition}

\begin{theorem}[Gauge local endpoint collapse to the canonical curvature-adjoint form]
\label{thm:gauge-local-endpoint-collapse}\lean{GaugeLane.endpointCollapse}
Assume the hypotheses of
Definition~\ref{def:gauge-local-endpoint-rigidity-package}. Then the forced
gauge-compatible endpoint class on the selected carrier has the canonical
curvature-adjoint representative
\[
\mathcal M_{\mathrm g}([A])=\delta_{\chi} dA=\delta_{\chi}F,
\]
where \(\delta_{\chi}\) denotes the adjoint of \(d\) with respect to the
canonical selected-cut pairing on the selected one-form carrier. Hence
\[
dF=0,
\qquad
\delta_{\chi}F=J+\mathsf R_{\mathrm g}.
\]
If the selected cut is lossless at leading order, so that
\(\mathsf R_{\mathrm g}=0\), this is the lossless intrinsic gauge law on the
selected one-form carrier.
\end{theorem}

(Proof in Appendix~\ref{sec:appendix-carrier}.)

\begin{corollary}[Recognizable Maxwell identity]
\label{cor:recognizable-maxwell-identity}\lean{GaugeLane.recognizableIdentity}
Assume the hypotheses of
Theorem~\ref{thm:gauge-local-endpoint-collapse}. Assume in addition that the
generated one-form carrier carries the canonical Hodge-side duality induced by
the chosen metric on \(M\), so the selected-cut adjoint \(\delta_{\chi}\) is
read as the codifferential \(\delta\).
Then
\[
dF=0,
\qquad
\delta F=J+\mathsf R_{\mathrm g}.
\]
If the selected cut is lossless at leading order, so that
\(\mathsf R_{\mathrm g}=0\), this is exactly the Maxwell curvature/source law,
with \(J\) recognized as the visible exchange current.
\end{corollary}

\begin{proof}
Substitute the Hodge-side identification \(\delta_{\chi}=\delta\) into
Theorem~\ref{thm:gauge-local-endpoint-collapse}.
\end{proof}

\begin{corollary}[Exact Schur form of the selected gauge remainder]
\label{cor:gauge-schur-remainder-identity}\lean{GaugeLane.schurRemainder}
Assume the hypotheses of
Theorem~\ref{thm:gauge-local-endpoint-collapse}. Assume in addition that the
same selected datum satisfies a canonical selected Schur bridge. Then there is
a unique gauge-projected Schur correction
\[
\Sigma_{h,\mathrm g}(z)\in\Omega^1(M)
\]
obtained by reading the common selected Schur correction \(\CCSigma{h}{z}\)
through the forced gauge endpoint identification on the selected carrier, and
\[
\mathsf R_{\mathrm g}=\Sigma_{h,\mathrm g}(z).
\]
Accordingly,
\[
dF=0,
\qquad
\delta_{\chi}F=J+\Sigma_{h,\mathrm g}(z).
\]
\end{corollary}

(Proof in Appendix~\ref{sec:appendix-carrier}.)

\subsubsection{Concrete Hidden-Channel Gauge Model}

The exact Schur form above becomes computationally informative once one
computes the returned gauge current on a concrete realized family. The
simplest solvable case is a transverse visible gauge mode locally coupled to
one hidden conserved current mode on the same primitive lattice. On that
model the gauge remainder is an exact current-side self-energy.

\begin{definition}[Coupled visible-shadow gauge mode model]
\label{def:coupled-visible-shadow-gauge-mode-model}\lean{GaugeLane.HiddenChannel.LatticeModel}
Assume the hypotheses of
Corollary~\ref{cor:gauge-schur-remainder-identity}. A
\emph{coupled visible-shadow gauge mode model} consists of:
\begin{enumerate}[label=(\roman*),itemsep=0.2em]
\item a one-dimensional primitive lattice with spacing \(\ell_{\chi}>0\)
      and a fixed transverse polarization one-form \(\varepsilon\);
\item a selected visible gauge mode and external visible current
      \[
      A_n(t)=a\,\varepsilon\,e^{i(nk\ell_{\chi}-\omega t)},
      \qquad
      J_n^{\mathrm{ext}}(t)=j_0\,\varepsilon\,e^{i(nk\ell_{\chi}-\omega t)};
      \]
\item a hidden transverse current mode
      \[
      B_n(t)=b\,\varepsilon\,e^{i(nk\ell_{\chi}-\omega t)};
      \]
\item visible and hidden spectral profiles
      \[
      \Omega_{\mathrm{g,vis}}(k)^2
      :=
      \frac{4c_{\chi}^2}{\ell_{\chi}^2}
      \sin^2\!\left(\frac{k\ell_{\chi}}{2}\right),
      \]
      \[
      \Omega_{\mathrm{g,hid}}(k)^2
      :=
      \nu_{\chi}^2
      +
      \frac{4\widetilde c_{\chi}^2}{\ell_{\chi}^2}
      \sin^2\!\left(\frac{k\ell_{\chi}}{2}\right),
      \]
      where \(c_{\chi}>0\), \(\widetilde c_{\chi}>0\), \(\nu_{\chi}\ge 0\),
      and \(\gamma_{\chi}\ge 0\);
\item a local coupling on the selected transverse sector such that the visible
      and hidden mode amplitudes satisfy
      \[
      \bigl(\Omega_{\mathrm{g,vis}}(k)^2-\omega^2\bigr)a
      =
      j_0+\gamma_{\chi}b,
      \]
      \[
      \bigl(\Omega_{\mathrm{g,hid}}(k)^2-\omega^2\bigr)b
      =
      \gamma_{\chi}a.
      \]
\end{enumerate}
\end{definition}

\begin{theorem}[Exact hidden-channel gauge current and branch splitting]
\label{thm:hidden-channel-gauge-current-and-branch-splitting}\lean{GaugeLane.HiddenChannel.splitting}
Assume the hypotheses of
Definition~\ref{def:coupled-visible-shadow-gauge-mode-model}. Then on every
mode with \(\Omega_{\mathrm{g,hid}}(k)^2\neq\omega^2\), exact elimination of
the hidden current amplitude gives the returned gauge current profile
\[
\sigma_{\chi}^{\mathrm{g,hid}}(k,\omega)
=
\frac{\gamma_{\chi}^2}{
\Omega_{\mathrm{g,hid}}(k)^2-\omega^2
}.
\]
Accordingly the visible response law is
\[
\bigl(\Omega_{\mathrm{g,vis}}(k)^2-\omega^2\bigr)a
=
j_0+\sigma_{\chi}^{\mathrm{g,hid}}(k,\omega)\,a,
\]
so the returned one-form remainder on that sector is
\[
\Sigma_{h,\mathrm g}^{\mathrm{hid}}(k,\omega)
=
\sigma_{\chi}^{\mathrm{g,hid}}(k,\omega)\,
a\,\varepsilon\,e^{i(nk\ell_{\chi}-\omega t)}.
\]
If \(j_0=0\), then
\[
(\omega^2-\Omega_{\mathrm{g,vis}}(k)^2)
(\omega^2-\Omega_{\mathrm{g,hid}}(k)^2)
=
\gamma_{\chi}^2,
\]
and the two exact mixed gauge branches are
\[
\omega_{\pm}(k)^2
=
\frac{
\Omega_{\mathrm{g,vis}}(k)^2+\Omega_{\mathrm{g,hid}}(k)^2
\pm
\sqrt{
\bigl(\Omega_{\mathrm{g,vis}}(k)^2-\Omega_{\mathrm{g,hid}}(k)^2\bigr)^2
+
4\gamma_{\chi}^2
}
}{2}.
\]
\end{theorem}

(Proof in Appendix~\ref{sec:appendix-carrier}.)

\begin{corollary}[Exact off-resonant gauge pole shifts]
\label{cor:exact-off-resonant-gauge-pole-shifts}\lean{GaugeLane.HiddenChannel.offResonantShifts}
Assume the hypotheses of
Theorem~\ref{thm:hidden-channel-gauge-current-and-branch-splitting} and
\(j_0=0\). Define
\[
\delta_{\chi}^{\mathrm g}(k)
:=
\Omega_{\mathrm{g,hid}}(k)^2-\Omega_{\mathrm{g,vis}}(k)^2.
\]
If \(\delta_{\chi}^{\mathrm g}(k)>0\), then
\[
\omega_-(k)^2
=
\Omega_{\mathrm{g,vis}}(k)^2
-
\frac{2\gamma_{\chi}^2}{
\delta_{\chi}^{\mathrm g}(k)
+
\sqrt{
\bigl(\delta_{\chi}^{\mathrm g}(k)\bigr)^2+4\gamma_{\chi}^2
}
},
\]
\[
\omega_+(k)^2
=
\Omega_{\mathrm{g,hid}}(k)^2
+
\frac{2\gamma_{\chi}^2}{
\delta_{\chi}^{\mathrm g}(k)
+
\sqrt{
\bigl(\delta_{\chi}^{\mathrm g}(k)\bigr)^2+4\gamma_{\chi}^2
}
}.
\]
In particular, away from hidden resonance the visible-like gauge pole is
shifted downward and the hidden current pole upward.
\end{corollary}

(Proof in Appendix~\ref{sec:appendix-carrier}.)

\begin{corollary}[Exact avoided crossing of the coupled gauge branches]
\label{cor:exact-avoided-crossing-coupled-gauge-branches}\lean{GaugeLane.HiddenChannel.avoidedCrossing}
Assume the hypotheses of
Theorem~\ref{thm:hidden-channel-gauge-current-and-branch-splitting} and
\(j_0=0\). Suppose there exists a selected mode \(k_\ast\) such that
\[
\Omega_{\mathrm{g,vis}}(k_\ast)^2
=
\Omega_{\mathrm{g,hid}}(k_\ast)^2
=:
\Omega_{\ast}^2.
\]
Then
\[
\omega_{\pm}(k_\ast)^2=\Omega_{\ast}^2\pm\gamma_{\chi},
\]
so the squared-frequency gap at the would-be crossing is exactly
\[
\omega_+(k_\ast)^2-\omega_-(k_\ast)^2=2\gamma_{\chi}.
\]
In particular, if \(\gamma_{\chi}>0\), the two gauge branches do not cross.
\end{corollary}

(Proof in Appendix~\ref{sec:appendix-carrier}.)

\begin{corollary}[Low-frequency induced exchange susceptibility]
\label{cor:low-frequency-induced-exchange-susceptibility}\lean{GaugeLane.HiddenChannel.lowFrequencySusceptibility}
Assume the hypotheses of
Theorem~\ref{thm:hidden-channel-gauge-current-and-branch-splitting} and
\(\nu_{\chi}>0\). Then as \((k,\omega)\to(0,0)\),
\[
\sigma_{\chi}^{\mathrm{g,hid}}(k,\omega)
=
\frac{\gamma_{\chi}^2}{\nu_{\chi}^2}
-
\frac{\gamma_{\chi}^2\widetilde c_{\chi}^2}{\nu_{\chi}^4}k^2
+
\frac{\gamma_{\chi}^2}{\nu_{\chi}^4}\omega^2
+
O(k^4+k^2\omega^2+\omega^4).
\]
Thus on the long-wave low-frequency sector the hidden return acts as an
induced local exchange susceptibility with static amplitude
\(\gamma_{\chi}^2/\nu_{\chi}^2\) together with the first dispersive
corrections.
\end{corollary}

(Proof in Appendix~\ref{sec:appendix-carrier}.)

The two-channel gauge model gives an exact realization of the returned gauge
current. On this model, the remainder manifests as gauge-pole shifts, avoided
crossing, and low-frequency induced exchange susceptibility on the visible
gauge sector.

\subsubsection{Exact Euclidean gauge shadow-flow}

The returned gauge correction admits an exact determinant-flow formulation
that is native to the tower calculus and structurally distinct from perturbative
running. This is the monograph's intrinsic exact gauge response law.

\begin{definition}[Selected gauge shadow-flow]
\label{def:selected-gauge-shadow-flow}\lean{GaugeLane.SelectedGaugeShadowFlow}
Assume the selected cut carries an intrinsic selected gauge carrier
$\mathcal Y_{h,\mathrm g}\subseteq \mathcal X_h$, and assume
$\mathcal Y_{h,\mathrm g}$ is finite-dimensional.
Let $P_{h,\mathrm g}:\mathcal X_h\to \mathcal Y_{h,\mathrm g}$ be the
orthogonal projection onto that carrier.

Let $T_h$ be the shadow-transport operator of
Definition~\ref{def:shadow-transport-invariants}. Define the
\emph{gauge shadow-transport compression} by
\[
T_{h,\mathrm g}
:=
P_{h,\mathrm g}\,T_h\,P_{h,\mathrm g}\big|_{\mathcal Y_{h,\mathrm g}}.
\]

Define the associated determinant flow and exact Euclidean gauge
shadow-flow by
\[
\Delta_{h,\mathrm g}(t):=\det(I+tT_{h,\mathrm g}),
\qquad
\tau_{h,\mathrm g}(t):=\frac{d}{dt}\log \Delta_{h,\mathrm g}(t),
\qquad t\ge 0.
\]
\end{definition}

\begin{theorem}[Exact Euclidean gauge shadow-flow law]
\label{thm:exact-euclidean-gauge-shadow-flow}\lean{GaugeLane.exactEuclideanGaugeShadowFlow}
Assume Definition~\ref{def:selected-gauge-shadow-flow}. Then:

\begin{enumerate}[label=\textup{(\roman*)}]
\item $T_{h,\mathrm g}$ is positive semidefinite.

\item If
$\Lambda_{h,\mathrm g}
=\{\lambda_{h,1},\dots,\lambda_{h,r_h}\}$
is the multiset of eigenvalues of $T_{h,\mathrm g}$ with multiplicity,
then
\[
\Delta_{h,\mathrm g}(t)=\prod_{j=1}^{r_h}(1+t\lambda_{h,j}),
\]
and
\[
\tau_{h,\mathrm g}(t)
=
\sum_{j=1}^{r_h}\frac{\lambda_{h,j}}{1+t\lambda_{h,j}}
=
\Tr\!\bigl(T_{h,\mathrm g}(I+tT_{h,\mathrm g})^{-1}\bigr).
\]

\item For every $n\ge 0$ and every $t\ge 0$,
\[
(-1)^n \tau_{h,\mathrm g}^{(n)}(t)
=
n!\,
\Tr\!\Bigl(
T_{h,\mathrm g}^{\,n+1}(I+tT_{h,\mathrm g})^{-(n+1)}
\Bigr)
\ge 0.
\]
Hence $\tau_{h,\mathrm g}$ is positive, decreasing, and completely
monotone on $[0,\infty)$.

\item If $\|tT_{h,\mathrm g}\|<1$, then
\[
\tau_{h,\mathrm g}(t)
=
\sum_{m=0}^{\infty}(-1)^m M_{m+1}^{(\mathrm g)}(h)\,t^m,
\]
where $M_k^{(\mathrm g)}(h):=\Tr(T_{h,\mathrm g}^k)$.

\item Let $E_{h,\mathrm g}$ be the spectral measure of $T_{h,\mathrm g}$,
and define the finite positive Borel measure
$\nu_{h,\mathrm g}(B)
:=\Tr\!\bigl(E_{h,\mathrm g}(B)\,T_{h,\mathrm g}\bigr)$.
Then
\[
\tau_{h,\mathrm g}(t)
=
\int_{[0,\infty)} \frac{1}{1+t\lambda}\,d\nu_{h,\mathrm g}(\lambda).
\]
In particular, $\tau_{h,\mathrm g}$ is a Stieltjes function.

\item $\tau_{h,\mathrm g}$ determines $\Delta_{h,\mathrm g}$ by
\[
\Delta_{h,\mathrm g}(t)
=
\exp\!\Bigl(\int_0^t \tau_{h,\mathrm g}(s)\,ds\Bigr),
\qquad
\Delta_{h,\mathrm g}(0)=1,
\]
and therefore determines the full spectral multiset
$\Lambda_{h,\mathrm g}$ and all moments $M_k^{(\mathrm g)}(h)$.
\end{enumerate}
\end{theorem}

\begin{proof}
Since $T_h$ is positive semidefinite by
Definition~\ref{def:shadow-transport-invariants}, its orthogonal
compression $T_{h,\mathrm g}=P_{h,\mathrm g}T_hP_{h,\mathrm g}$ is also
positive semidefinite: for every $u\in\mathcal Y_{h,\mathrm g}$,
$\langle u,T_{h,\mathrm g}u\rangle
=\langle P_{h,\mathrm g}u,T_hP_{h,\mathrm g}u\rangle
=\langle u,T_hu\rangle\ge 0$.
This proves~(i).

Because $\mathcal Y_{h,\mathrm g}$ is finite-dimensional and
$T_{h,\mathrm g}$ is self-adjoint, there is an orthonormal basis in which
$T_{h,\mathrm g}$ is diagonal with nonneg\-ative entries
$\lambda_{h,1},\dots,\lambda_{h,r_h}$. Therefore
$\Delta_{h,\mathrm g}(t)=\prod_j(1+t\lambda_{h,j})$.
Differentiating $\log\Delta_{h,\mathrm g}$ gives
$\tau_{h,\mathrm g}(t)=\sum_j\lambda_{h,j}/(1+t\lambda_{h,j})$.
The trace identity follows from the diagonal form of
$T_{h,\mathrm g}(I+tT_{h,\mathrm g})^{-1}$.
This proves~(ii).

Termwise differentiation yields
$\tau_{h,\mathrm g}^{(n)}(t)
=(-1)^n n!\sum_j\lambda_{h,j}^{n+1}/(1+t\lambda_{h,j})^{n+1}$,
which equals
$(-1)^n n!\,\Tr(T_{h,\mathrm g}^{n+1}(I+tT_{h,\mathrm g})^{-(n+1)})$.
Each summand is nonneg\-ative since $\lambda_{h,j}\ge 0$ and
$1+t\lambda_{h,j}>0$. Complete monotonicity follows.
This proves~(iii).

The geometric expansion
$\lambda_j/(1+t\lambda_j)=\sum_{m\ge 0}(-1)^m\lambda_j^{m+1}t^m$
converges absolutely for $|t\lambda_j|<1$. Summing over~$j$ gives the
moment series with $M_k^{(\mathrm g)}(h)=\sum_j\lambda_{h,j}^k$.
This proves~(iv).

By functional calculus,
$T_{h,\mathrm g}(I+tT_{h,\mathrm g})^{-1}
=\int\frac{\lambda}{1+t\lambda}\,dE_{h,\mathrm g}(\lambda)$.
Taking traces gives
$\tau_{h,\mathrm g}(t)=\int(1+t\lambda)^{-1}\,d\nu_{h,\mathrm g}(\lambda)$
with $\nu_{h,\mathrm g}$ positive (since $E_{h,\mathrm g}$ and
$T_{h,\mathrm g}$ are positive). This proves~(v).

Part~(vi) follows from $\Delta_{h,\mathrm g}(t)
=\exp(\int_0^t\tau_{h,\mathrm g}(s)\,ds)$ and the fact that the
characteristic polynomial $\prod_j(1+t\lambda_{h,j})$ determines the
eigenvalue multiset $\Lambda_{h,\mathrm g}$ and hence all moments.
\end{proof}

\begin{corollary}[Finite-horizon nonlogarithmic distinctness]
\label{cor:finite-horizon-nonlogarithmic-distinctness}\lean{GaugeLane.finiteHorizonNonlogarithmicDistinctness}
At every fixed horizon~$h$, the exact gauge shadow-flow
$\tau_{h,\mathrm g}(t)$ is a rational Stieltjes function of~$t$,
with simple poles only at the negative real numbers
$-\lambda_{h,j}^{-1}$ corresponding to the nonzero eigenvalues of
$T_{h,\mathrm g}$. In particular, the monograph-native exact law is not
logarithmic at finite resolution. Any logarithmic running can arise only
after passage to a continuum spectral limit of the measures
$\nu_{h,\mathrm g}$.
\end{corollary}

\begin{proof}
By Theorem~\ref{thm:exact-euclidean-gauge-shadow-flow}(ii),
$\tau_{h,\mathrm g}(t)=\sum_j\lambda_{h,j}/(1+t\lambda_{h,j})$
is rational with the stated poles. A rational function is not logarithmic.
\end{proof}

\begin{theorem}[Spectral-limit gauge shadow-flow]
\label{thm:spectral-limit-gauge-shadow-flow}\lean{GaugeLane.spectralLimitGaugeShadowFlow}
Assume a sequence of horizons $h_n$ such that the finite positive measures
$\nu_{h_n,\mathrm g}$ of
Theorem~\ref{thm:exact-euclidean-gauge-shadow-flow}(v) converge weakly to
a finite positive Borel measure $\nu_{\mathrm g}$ on $[0,\infty)$.
Define
\[
\tau_{\mathrm g}(t)
:=
\int_{[0,\infty)}\frac{1}{1+t\lambda}\,d\nu_{\mathrm g}(\lambda),
\qquad t\ge 0.
\]
Then for every $t\ge 0$,
$\tau_{h_n,\mathrm g}(t)\to \tau_{\mathrm g}(t)$.
\end{theorem}

\begin{proof}
For each fixed $t\ge 0$, the kernel $f_t(\lambda)=(1+t\lambda)^{-1}$
is bounded and continuous on $[0,\infty)$. By weak convergence,
$\int f_t\,d\nu_{h_n,\mathrm g}\to\int f_t\,d\nu_{\mathrm g}$.
By Theorem~\ref{thm:exact-euclidean-gauge-shadow-flow}(v), the left side
is $\tau_{h_n,\mathrm g}(t)$ and the right is $\tau_{\mathrm g}(t)$.
\end{proof}
